% TOP_PROBLEM: {"id": 5300003, "title": "Continuity of polynomial mating", "category_id": 11, "category": {"id": 11, "name": "dynamical_systems", "display_name": "Dynamical Systems", "description": "Problems about long-term behavior of deterministic systems, Hamiltonian dynamics, and periodic orbits.", "slug": "dynamical-systems", "order_index": 11, "created_at": "2026-07-31T15:26:25.671Z"}, "status": "partially_solved"}
% FIRST_POSTED: null
% ENTRY_KIND: statement_only
% Imported from: batch04.tex; UnsolvedMath ID 5300003
\documentclass[11pt]{article}
\usepackage{fontspec}
\setmainfont{Latin Modern Roman}
\setsansfont{Latin Modern Sans}
\usepackage{amsmath,amssymb,mathtools}
\usepackage{unicode-math}
\setmathfont{Latin Modern Math}
\usepackage{xurl}
\usepackage[hidelinks]{hyperref}
\newcommand{\sref}[2]{\hyperlink{source:#1:#2}{[S#2]}}
\newcommand{\cataloguescope}[1]{\par\noindent\textbf{Question.} #1\par}
\begin{document}
% UnsolvedMath ID 5300003; source catalogue code AMR-052-0003
\setcounter{section}{5300002}
\section{Continuity of polynomial mating}\label{5300003}
{\small\sffamily UnsolvedMath ID 5300003 · Dynamical Systems}
\subsection{Definitions and mathematical statement}

\paragraph{Shared notation.}$\mathbb N=\{1,2,\ldots\}$, and $\mathbb Z,\mathbb Q,\mathbb R,\mathbb C$ denote the integers, rationals, reals and complex numbers. Any other conventions are specified locally.


For a complex polynomial $P$ of degree at least two, its filled Julia set $K(P)$ is the set of points with bounded forward orbit, and $J(P)=\partial K(P)$ is its Julia set. An external ray is a radial curve in the coordinate near infinity that conjugates $P$ to its leading power map, continued into the basin of infinity. Its angle is measured in $\mathbb R/\mathbb Z$. A critical point is a zero of $P'$. Topological conjugacy means a homeomorphism $h$ satisfying $h\circ f=g\circ h$. A quasiconformal homeomorphism has bounded infinitesimal eccentricity; quasiconformal surgery cuts and glues dynamical planes and finds an invariant complex structure that can be straightened to a holomorphic map.
For monic polynomials $P,Q$ of the same degree $d\geq2$, compactify each dynamical plane by a circle at infinity. On that circle the dynamics is $\theta\mapsto d\theta$. Glue the circles by identifying angle $\theta$ for $P$ with angle $-\theta$ for $Q$, and then collapse the external rays to points. The resulting quotient and induced map, when well defined, are the topological mating. The construction itself can fail to give a sphere or a rationally realizable map.
In questions about variation, polynomial coefficients carry their usual topology; on a fixed-degree space this is equivalent to locally uniform convergence. Rational maps are considered up to conformal conjugacy, with the topology inherited from their fixed-degree coefficient space. The identifications used in the surgery are kept as part of the construction when the inputs vary.
\paragraph{Statement recorded in the supplied source.}
Where mating has a rational realization, is that realization continuous when one input polynomial is fixed and the other varies continuously? Is mating jointly continuous as a function of both input polynomials? \sref{5300003}{2}
\subsection{Short English statement}
Does mating depend continuously on each varying input, and is it a jointly continuous operation on pairs of polynomials?
\subsection{Sources}
\begin{description}
\item[\hypertarget{source:5300003:1}{S1}] ULAM.AI and the UnsolvedMath Contributors, UnsolvedMath: A Curated Collection of Open Mathematics Problems, record 5300003 (AMR-052-0003). CC BY 4.0. \url{https://huggingface.co/datasets/ulamai/UnsolvedMath}.
\item[\hypertarget{source:5300003:2}{S2}] Ben Bielefeld, Questions in Quasiconformal Surgery, in Ben Bielefeld and Mikhail Lyubich (editors), Problems in holomorphic dynamics (1992). Question 3, p. 3. \url{https://arxiv.org/pdf/math/9205209}.
\end{description}
\label{end:5300003}
\end{document}
